% 根拠：E2（条件）、E3（現行結果）、E4（比較手法）、E5/E6/E7（構成要素・感度）。
% 執筆時の注意：旧Observed-lag AR/Raw × BF/GLRTと現行方式の比較を混ぜない。
% 入力列処理・標準化等の条件差を照合する。RUNの正式結果は前回確認時に未完了。
\section{評価実験}
\label{sec:experiment}

\subsection{どんなデータで，何を正解として調べるか}
\begin{itemize}
    \item RCAEvalは，障害原因分析の方法を比較するための公開データ集．システムに意図的な不具合を起こし，各所の観測値と，原因サービス・開始時刻を記録している．そのうち三つのシステム，計375件の障害事例を使う．
    \item 障害前後からそれぞれ最大600時点を取り出す．正解は不具合を起こしたサービスとし，どの方法にも同じ障害事例を与える．データ版などの再現に必要な情報は，後の稿で条件表へまとめる．
    \item 注意：使える開始時刻や学習に使う区間なども比較条件である．既存方法を今回のデータに合わせて変更した部分は，後で明記する．
\end{itemize}

\subsection{問い1：原因をどの程度上位に挙げられるか}
\begin{itemize}
    \item 提案手法と既存方法について，「第一候補が正解だった割合」「上位三つ・五つに正解が入った割合」を比較する．この上位$k$件の正解率をAC@$k$と呼ぶ．
    \item 表案：方法ごとの正解率を，三つのシステム別と全体で示す．候補を出せなかった事例も採点に含める．計算にかかる時間は，測定条件をそろえて補足する．
    \item 注意：まだ計算できていない結果は空欄として残し，得られた範囲だけで比較する．数値は保存済みの実験結果を照合してから記入する．
\end{itemize}

\subsection{問い2：どの工夫が役立っているか}
\begin{itemize}
    \item 時系列予測を使わず観測値の分布を直接比べる場合，障害後も実測値を使って予測する場合，予測の不確かさを補正しない場合を，提案方法と比較する．
    \item ベイズ的な比較の役割は，BFの代わりに一般化尤度比を使う比較から調べる．変えた要素と固定した条件を表にし，複数の変更による差を一つの工夫の効果としない．
    \item 図表案：同じ障害事例で正解の順位が上がった件数・下がった件数，第一候補への出入りを示す．TODO：最終方式について，どの比較まで同じ条件で確認できているかを整理する．
\end{itemize}

\subsection{問い3：結果はどの条件で変わるか}
\begin{itemize}
    \item AR次数などの設定や，秒からミリ秒への単位変換によって順位が変わるかを調べる．不具合の種類ごとの得意・不得意も見る．
    \item 正常時のデータを前半・後半に分け，後半を仮の「障害後」として得点を計算する．不具合がなくても高得点になるなら，その得点を障害の確かな証拠と読めるかを考える．
    \item 注意：同じデータを何度も見て設定を選んでいる点を明記する．小さな精度差を最適性の証明とせず，統計的な不確かさも併記する．
\end{itemize}
